\documentclass[11pt]{article}

\newif\ifanon
\anonfalse

\usepackage[margin=1in]{geometry}
\usepackage[utf8]{inputenc}
\usepackage[T1]{fontenc}
\usepackage{times}
\usepackage{amsmath,amssymb}
\usepackage{booktabs}
\usepackage{tabularx}
\usepackage[round]{natbib}
\usepackage{xcolor}
\usepackage{enumitem}
\usepackage{microtype}
\usepackage{titling}
\usepackage{float}
\usepackage{pgfplots}
\pgfplotsset{compat=1.18}
\usepackage{hyperref}

\setlength{\droptitle}{-1cm}
\pretitle{\begin{center}\Large\bfseries}
\posttitle{\end{center}\vspace{0.3cm}}

\hypersetup{
  colorlinks=true,
  linkcolor=black,
  citecolor=blue!50!black,
  urlcolor=blue!50!black
}
\ifanon\else
\hypersetup{
  pdftitle={Does the Direction Track the Process? A Preregistered Validity Battery for Published Hidden-State Directions},
  pdfauthor={Sam Ramdan}
}
\fi
\setlist{itemsep=2pt,topsep=3pt}
\setlength{\belowcaptionskip}{6pt}

\definecolor{kdec}{HTML}{2A78D6}
\definecolor{kans}{HTML}{EB6834}
\definecolor{kcomp}{HTML}{1BAF7A}

\newcommand{\dhat}{\hat{d}}
\newcommand{\verdict}[1]{\textsc{\MakeLowercase{#1}}}
\newcommand{\aidisclosure}{An AI system (Claude, Anthropic), working under the author's direction, reviewed the literature, wrote the preregistrations and code, ran the experiments and analyses, and drafted this paper. The full statement is in the paragraph ``Use of AI assistance'' before the references.}

\title{Does the Direction Track the Process?\\ A Preregistered Validity Battery for Published Hidden-State Directions}

\ifanon
\author{Anonymous authors\thanks{\aidisclosure}}
\else
\author{
  Sam Ramdan\thanks{\aidisclosure}\\
  Independent Researcher\\
  \texttt{sr4729820@gmail.com}\\
  ORCID \href{https://orcid.org/0009-0004-7723-4240}{0009-0004-7723-4240}
}
\fi
\date{October 2026}

\begin{document}

\maketitle

\begin{abstract}
A linear direction in a language model's hidden states is usually defended by two facts: it separates prompts that differ by design, and adding it changes the output as expected. Neither shows that it follows the process it is named after. We preregistered a battery of checks and apply it to nine published directions on five open 7-8B models (six configurations: Qwen3-8B with thinking off and on). The battery asks whether a direction tracks a per-item target measured from the model's own behaviour, survives stimuli matched on template, length and wording, beats two label-free directions built from the same calibration prompts, steers beyond random pushes matched on output KL divergence, and acts through later layers rather than as an output bias. An oracle direction, the difference of means between items high and low on the target, asks whether a push of one calibration gap can move the target. Both controls pass: the refusal and truth directions are valid drivers. None of six process directions receives a valid verdict. Two do not reproduce their authors' headline steering effect; a third, LiReFs, passes only a reproduction gate moved to class separation after its pilot. The oracle reads the process in 17 of 19 process runs, yet its push passes the steering check in none of them, against all six control runs. This null concerns the oracle's push, not the target: in at least three runs the published direction pushes where the oracle does not. In exploratory pushes up to eight gaps, answer choices moved along the oracle in 5 of 9 runs, settling depth and chain-of-thought necessity in 1 of 8, and the published directions moved their own targets in 9 of 19 process runs. Weight-randomised models reproduce calibration separation for four of nine directions, and norm-matched random baselines overstate steering effects against KL-matched ones, mostly in two of 23 pairs.
\end{abstract}

\section{Introduction}

Interpretability papers often claim that a linear direction $d$ in a model's residual stream \emph{tracks} a process (its projection rises and falls with the process) or \emph{drives} it (adding $d$ changes the process). The usual evidence has two parts. The direction separates two sets of prompts that differ in the process by construction, and adding it at some norm changes the output in the expected way. Both parts are about the prompts and the output, not about the process.

Sanity checks for this kind of evidence have a history. Saliency maps that look plausible can be reproduced by a network with random weights \citep{adebayo2018sanity}, and probe accuracy has to be read against what a probe learns from arbitrary labels \citep{hewitt2019control}. For directions, several 2026 studies find that a probe direction can read a target that pushing along it does not move \citep{xun2026lagged}; Section~\ref{sec:related} reviews them. A worked case motivated the present battery. In earlier work, a compute--retrieve (CR) axis separated computation prompts from retrieval prompts at AUC $\ge 0.988$ in twelve models \citep{ramdan2026crv1}, reversed in three of five base models under exact format and length matching \citep{ramdan2026fingerprint}, and in a preregistered causal test neither tracked nor moved the depth at which an arithmetic answer settles; it moved reply format instead, as did a direction built from the same sixty calibration prompts without their labels (long minus short within each class), which also tracked settling depth better \citep{ramdan2026crv2}. Reports of injected thoughts in Qwen3 models likewise ran mainly through one content-free direction that acted as a general yes-bias \citep{ramdan2026injected}. None of this was visible from calibration AUC or from a fixed-norm steering demonstration.

The checks involved are generic and cheap. Measure the process per item from what the model does, not from the prompt's class. Match stimuli on template, length and vocabulary. Compare with a label-free direction from the same calibration prompts. Compare with random pushes matched on the change they cause in the output distribution rather than on norm. Find the stage at which the push's effect is decided. Each check has precedents taken one at a time, KL-matched random controls among them, and the gap between reading a variable and pushing it, which dominates our results, has been reported repeatedly in 2026 (Section~\ref{sec:related}). What we add is narrower: process targets measured per item from what the model does; an oracle direction built from that target, as a check on whether a failure of the published direction says anything about the direction, within limits we set out in Section~\ref{sec:oracle}; label-free competitors built from the same calibration prompts; calibration on weight-randomised models; and preregistered breadth, with one fixed battery applied to nine published directions on five open 7--8B models, checked against a positive control (passed) and a known-answer case (partly retrodicted). This is the analogue, for directions, of the sanity checks \citet{adebayo2018sanity} proposed for saliency maps.

We preregistered the battery, its pass criteria, its verdict rules and seven hypotheses before writing any code, and froze per-direction settings after pilots on separate items. We apply it to nine published directions on five open checkpoints of 7 to 8B parameters (six configurations, counting Qwen3-8B with thinking off and on): two controls (refusal, \citealp{arditi2024refusal}; truth, \citealp{marks2024geometry}), the CR axis as a known-answer case, and six directions that claim to track or drive a process.

The controls pass, so the battery is not built to fail: refusal and truth receive \verdict{VALID DRIVER} overall. None of the six process directions does (H5). Two are \verdict{NOT REPRODUCED} overall: our difference-of-means rebuild of known-entity awareness on ungated chat models does not produce the refusal effect reported with sparse-autoencoder latents in Gemma-2, and reasoning strength reproduces only in its home model. A third, LiReFs, passes a reproduction gate that was set on class separation after its pilot; its intervention never improved held-out reasoning, and under the gate as Stage~1 defined it (the authors' headline effect) it too would be \verdict{NOT REPRODUCED}. The commonest verdict is \verdict{TARGET NOT LINEAR}: an oracle direction built from the process reads it in 17 of 19 process runs, yet its one-gap push passes the steering check in none of them, against all six control runs. That is a null for the oracle's push, not evidence that the target cannot be pushed: the oracle is a reading direction, and in at least three runs the published direction moves its target where the oracle does not (Section~\ref{sec:oracle}). Exploratory pushes of up to eight gaps separate the targets only in part. Answer choices (hallucinating, giving in to the user, following a perturbed step) moved along the oracle in 5 of 9 runs; settling depth and chain-of-thought necessity moved in 1 of the 8 runs evaluable at four or eight gaps; backtracking was pushed along the oracle only to two gaps and did not move, while the published backtracking vector moves its keyword rate at its authors' strength of 7.7 gaps. The published directions moved their own targets in 9 of the 19 process runs at some strength, but LiReFs never moved settling depth or chain-of-thought necessity and the sycophancy vector never moved giving in. The only \verdict{NUISANCE} verdict falls on reasoning strength in its home model, where the rebuilt vector is nearly identical to the label-free first principal component and does no better than it. Behind the verdicts, known-entity awareness tracks its process and beats both label-free competitors in all three models, and true thinking tracks in all three and beats both in two; a second reasoning direction, like the CR axis, moves reply format instead of computation; and the sycophancy persona vector is higher when the user is right than when the user is wrong. Weight-randomised models reproduce calibration separation for four of nine directions, and norm-matched random baselines overstate steering effects relative to KL-matched ones (H4), although two of 23 pairs carry most of the excess.

\section{Related work}\label{sec:related}

\paragraph{Readable but not pushable.} The dissociation that dominates our results, a target that a linear direction reads well but cannot push, has been reported repeatedly in 2026, and we do not claim it as a finding of our own. The closest study is \citet{xun2026lagged}. Across the Pythia suite (160M to 12B parameters, eight training checkpoints, four answer-format task families), a logistic probe reads the target from the first checkpoint on, yet adding the probe direction at the probed site stays within a norm- and site-matched random null in 43 of 48 model--checkpoint cells, with flat dose--response. The protocol was preregistered, and a two-size OLMo-2 replication kept the direction of the findings at weaker magnitude. Their probe direction plays the role of our oracle. Other settings show the same gap. A direction that detects fake entities perfectly lies about $83^\circ$ from the direction that produces refusal in Gemma~2-2B-it, and $78^\circ$ to $83^\circ$ from it across four models \citep{galeone2026detection}. Across seven models from 117M to 7B parameters, linear probes detect hallucination above chance in the larger ones, while steering along the probe direction corrects hallucinations in none of the seven \citep{roy2026detection}. A hazard probe at 98\% AUROC in Qwen2.5-7B-Instruct coexists with steering that leaves most missed hazards uncorrected \citep{basu2026actionability}. The SAE features a deployed truth probe weights most overlap little with those that drive the model's answer \citep{tiwari2026decodability}. Reasoning-correctness directions in a small trained model are decodable but do not improve reasoning under intervention \citep{nwadike2026correctness}. In agent models, a belief that a request is unambiguous is decodable, causally inert and orthogonal to the decision to ask about it \citep{zavattari2026interventions}. What our study adds to this line is the oracle used as a diagnostic inside a battery for published directions: positive controls on which the same construction does push (refusal and truth), per-item process targets for six published process claims, and a comparison across kinds of target (Section~\ref{sec:strength}). A null push of the oracle is informative only about that push (Section~\ref{sec:oracle}).

\paragraph{Sanity checks and control tasks.} \citet{adebayo2018sanity} showed that several saliency methods give nearly the same maps for a trained network and for one with randomised weights. \citet{hewitt2019control} proposed control tasks so that probe accuracy is read against what a probe can learn from arbitrary labels. \citet{korznikov2026sae} recently ran random baselines against sparse autoencoders, and \citet{kutluozen2026linear} show for vision models that any sufficiently rich feature map, a random one included, admits a linear readout, so probe accuracy alone does not establish a linear representation. Our Check~0b is the model-randomisation test for directions, and our label-permuted directions are the control-task analogue.

\paragraph{Process-defined labels.} Several of the directions we test were built from labels defined by behaviour: known versus unknown entities by whether the model answers correctly \citep{ferrando2025entities}, annotated reasoning behaviours \citep{ward2025backtracking,venhoff2025steering}, and the True Thinking Score \citep{zhao2025truethinking}. \citet{kirch2026offpolicy} show that probes trained on what the monitored model actually did generalise better than probes trained on off-policy data. None of these re-scores other authors' directions against per-item process targets, which is our Check~A.

\paragraph{Matched stimuli and confounds.} Matched contrasts are standard in causal-abstraction benchmarks and in \citet{marks2024geometry}'s curated sets. \citet{sahoo2026format} found that probes for reasoning modes detect task format and fall to chance once format variables are regressed out; \citet{devbunova2026format} found the same for evaluation awareness; \citet{sadhu2026compliance} found an activation readout matched by a bag-of-words model, and \citet{ghosh2026confidence} a prompt-only TF-IDF classifier that matches confidence probes in most retrieval settings. \citet{petrov2026topic} showed that refusal directions \emph{built} from topic-matched contrasts are too small to steer, which is why we use matched sets only to evaluate directions.

\paragraph{Competitor directions.} Check~C compares a direction with directions built from the same calibration prompts without their labels. Two recent checks are close in spirit. \citet{xu2026competing} compare a steering vector with the vector built the same way for a competing target, and find the vectors for two ethical frameworks interchangeable. \citet{subhnil2026failures} find that no hidden-state probe locates the failing step of a reasoning trace better than a rule based on trace length alone. Our competitors need no second label: they ask whether the label adds anything beyond the dominant variation of the prompts the direction was built from.

\paragraph{Random-direction controls.} Many steering papers compare with random directions of equal norm. Equal norm can be a weak control: in \citet{li2023iti}'s Table~3 a random direction changed the output distribution by a KL of 0.02 against 0.27 for the mass-mean shift at the same strength. At a given perturbation size, contrastive and SAE directions, among others, change a model's downstream output more than isotropic random directions do \citep{ferreiradasilva2026error}, which is one reason an equal-norm reference is weak. SteerCheck \citep{luo2026steercheck} matches off-target KL in its preregistered audit and finds that sign-randomised null directions built the same way as the steering direction often keep substantial alignment with it (alignment leakage); \citet{ferrara2026introspection} matches random controls on Jensen--Shannon divergence. We use KL matching as a standard and do not claim it as new.

\paragraph{Where steering acts.} \citet{gao2026steering} show that common steering methods follow the answer-extraction index rather than the semantic label, through a low-rank output-sensitive component. \citet{billa2026lap} predict where steering succeeds from logit-lens accessibility: trained probes separate concepts at every layer, but difference-of-means steering works where the concept is aligned with the unembedding, usually late, and not at the conventional middle layer. In three omni-modal models, the layer at which a probe reads best is not the layer at which steering works, which falls in a narrow mid-to-late range of depth \citep{wang2026readbest}. \citet{yu2026answerbasin} argue that behaviour-level probing and steering directions are organised by the model's own distribution over answers rather than by the concept they are named for. \citet{yuan2026specificity} compare late-layer steering with a direct logit bias. Our Check~E2 is a fixed version of that comparison.

\paragraph{Multi-claim audits.} \citet{wu2026steerable} audit four claim types over seventeen models and show that a scaling trend in steerability is an artefact of units and operating points. Their audit concerns static concepts; ours concerns process claims, adds per-item process targets and structured label-free competitors, and includes positive controls and a known-answer case that test the battery itself.

\section{The battery}\label{sec:battery}

\subsection{Definitions}
A \emph{direction claim} is a tuple (model, layer $\ell^*$, reading position, vector $d$, process $X$, kind $\in$ \{tracks, drives\}). The \emph{calibration set} is the labelled prompts the authors used to build $d$, used unchanged. The \emph{process target} $T$ is a per-item number or binary outcome measured from the model's own behaviour or internals on that item (whether it refused, whether its answer settled late, whether its next sentence backtracks); it is never the label of the prompt's class. The \emph{projection} is $P = h_{\ell^*}(x)\cdot\dhat$ at the authors' reading position. The \emph{gap} $g$ is the mean projection of the positive calibration class minus that of the negative class; it is the unit of push size. A \emph{push} adds $\alpha\dhat$ to the output of block $\ell^*$ at every position from the reading position on, including generated tokens. Its \emph{size} is the \emph{summed KL}: the KL divergence from the clean to the pushed next-token distribution, summed over every position of the clean continuation under teacher forcing. Settling depth is expressed as a fraction of the model's depth; where it helps we also give layers.

\subsection{Control directions}
For every (direction, model) pair we build, from the direction's own calibration set:
\begin{itemize}
\item $N_\text{len}$, the within-class length direction: in each class, the mean activation of the longer half of the calibration prompts minus that of the shorter half, averaged over classes;
\item $N_\text{pc1}$, the first principal component of the calibration activations after centring each class on its own mean, signed to project positively on $d$;
\item $N_\text{perm}$, twenty label-permuted difference-of-means directions;
\item the oracle $o$, the difference of means between the top and bottom terciles of $T$ on an oracle split of 120 items disjoint from every test item;
\item twelve Gaussian random directions per model, pushed over a norm grid that brackets the KL of every tested push.
\end{itemize}
$N_\text{len}$ and $N_\text{pc1}$ are the label-free competitors of Check~C. The oracle asks whether the target can be read and pushed linearly at $\ell^*$. It is built to read the target, and nothing optimises it to push, so a push along it that moves the target shows that the target is pushable at $\ell^*$, while a null push does not show the converse. Stage~1 treats an oracle failure as making the direction's own failure uninformative; we keep that rule for the verdicts and examine what it can and cannot show in Section~\ref{sec:oracle}.

\subsection{Checks}
A check is \emph{passed} when the 95\% CI excludes zero in the claimed sign and the point estimate reaches the smallest effect size of interest (SESOI). It is \emph{failed} when the CI bound on the claimed side falls short of the SESOI. Otherwise it is inconclusive. CIs come from an item bootstrap with 2,000 resamples; Check~D uses a two-level bootstrap over items and random directions.

\begin{description}[leftmargin=1.2em,style=nextline]
\item[Check 0 (reproduction and randomised-model calibration).] 0a rebuilds the direction with the authors' code, data, layer rule and reading position, and reproduces their headline effect within a tolerance fixed after the pilot; a direction that fails is \verdict{NOT REPRODUCED} and leaves the test. 0b applies the same recipe to a copy of the model with re-initialised weights and reports how well that direction separates held-out calibration prompts (descriptive; it enters only H3).
\item[Check A (tracking).] Spearman $\rho(P,T)$, or the AUC of $P$ for a binary $T$, on a process test set in one template and answer format. SESOI $\rho = 0.20$ or AUC $= 0.60$. If $T$ correlates with prompt length ($|\rho| > 0.1$), the partial statistic is primary.
\item[Check B (label survives matching).] The AUC of $P$ for the class label on a set matched for template, token length and vocabulary, accepted only if a bag-of-tokens classifier separates it at AUC $\le 0.65$. Pass: AUC $\ge 0.70$ with CI lower bound above 0.5. B gives a separate \emph{label} verdict.
\item[Check C (beats label-free directions).] Paired bootstrap differences $\rho_d - \rho_N$ (tracking) and $\text{effect}_d - \text{effect}_N$ at matched KL (steering), for $N \in \{N_\text{len}, N_\text{pc1}\}$. Pass: $d$ beats both. A \verdict{NUISANCE} flag is raised when a nuisance direction on its own meets the pass criterion of A or D and $d$ does not beat it.
\item[Check D (steers beyond matched random pushes).] A push of $\pm 1\,g$, each sign on the items where it can move $T$. For each item and random direction, the change in $T$ is interpolated linearly in log KL between the two random norms that bracket the push's KL; the effect is the mean over items of the push's change minus the random reference. Pushes whose KL falls outside the grid are excluded and counted: at one gap, against the Gaussian reference, this excluded 1 of about 6,000 item pushes of the published directions and 28 of the oracle's, 22 of them in one cell. The same effect is computed against the label-permuted directions. A sign passes when, against both references, the CI excludes zero in the claimed direction and the estimate reaches the target's SESOI. Check~D passes if at least one sign passes and the other sign is not a significant reversal. The oracle's Check~D is read on the sign on which $d$ passes, if any; otherwise it passes if either sign passes. The effect against random directions of the push's norm, without KL matching, is reported for H4.
\item[Check E (effect handed on).] Run only for pushes that pass D. E1 re-imposes the clean component along $\dhat$ at each later block and finds the hand-off layer, the first at which at least half the effect survives. E2 adds the push only at the last block, at the same KL. Pass: hand-off before the last quarter of the remaining depth and E2 recovering less than half the effect. Otherwise the push is flagged \verdict{OUTPUT-BIAS}.
\end{description}

\subsection{Verdicts}
Each (direction, model) pair gets the first verdict that applies, in Stage~1's order: \verdict{NOT REPRODUCED} (0a fails); \verdict{TARGET NOT LINEAR} (the oracle fails A, or, for a drive claim, D); \verdict{NUISANCE} (the flag of Check~C); \verdict{FORMAT-MEDIATED} (it passes D with the reply format free but not with the format fixed by a prefill, Stage~1's format clamp) or \verdict{OUTPUT-BIAS} (the flag of Check~E); \verdict{VALID TRACKER} (it passes A and C-tracking, and D is not passed or not claimed); \verdict{VALID DRIVER} (it passes every check its claim requires: A and C-tracking, and D, C-steering and E); \verdict{NOT SPECIFIC} (it passes A or D, but C is inconclusive); \verdict{NOT SHOWN} (anything else). No target here was measured both with the format free and under a prefill (settling depth, for example, is measured only under the prefill), so \verdict{FORMAT-MEDIATED} was never evaluated. We keep the preregistered verdict names, but \verdict{TARGET NOT LINEAR} establishes less than its name: that the oracle did not read the target, or that its one-gap push did not pass Check~D, not that no linear push at $\ell^*$ moves the target (Section~\ref{sec:oracle}). The analysis code counts an oracle as failing whenever it does not pass; in every \verdict{TARGET NOT LINEAR} verdict reported here its one-gap push failed or reversed in both signs against at least one reference, so this choice decides none of them.

Check~B gives the label verdict: \verdict{LABEL SURVIVES MATCHING} (B passes), \verdict{REVERSED} (B's CI lies entirely below 0.5: on matched stimuli the projection separates the classes in the direction opposite to the label) or \verdict{LABEL DOES NOT SURVIVE}. A direction's overall verdict is the one it receives in at least two of its three panel models, otherwise \verdict{MIXED}. No correction is made across models or directions; all CIs are reported.

\subsection{Hypotheses}
\begin{description}[leftmargin=1.2em,style=nextline]
\item[H1 (positive control).] Refusal receives \verdict{VALID DRIVER} in at least two of three models. If it fails, the failing checks are declared miscalibrated and verdicts are recomputed without them.
\item[H2 (negative control).] Label-permuted directions fail A and D in every model. Stage~1 calls this an implementation check whose failure ``means a bug, not a finding''; Section~\ref{sec:discussion} explains why we do not read the outcome that way.
\item[H3 (Adebayo analogue).] For at least half of the directions, the recipe applied to a weight-randomised model separates held-out calibration prompts at AUC $\ge 0.80$ (a direction counts if this holds in at least two of its three models).
\item[H4.] Effects measured against norm-matched random directions exceed effects against KL-matched ones (mean over reproduced pairs of the difference in SESOI units, bootstrap CI above zero).
\item[H5.] Fewer than half of the process directions receive an overall \verdict{VALID TRACKER} or \verdict{VALID DRIVER}.
\item[H6.] At least two process directions receive \verdict{NUISANCE} in at least one model.
\item[H7 (retrodiction).] The battery reproduces the conclusions of \citet{ramdan2026crv2} for the CR axis in Qwen2.5-7B-Instruct and Qwen3-8B: A fails for settling depth, D stays below one layer, the verdict is \verdict{NUISANCE} or \verdict{NOT SHOWN}, and on shown working the axis does not beat $N_\text{len}$.
\end{description}

\section{Directions, models and procedure}\label{sec:directions}

\begin{table}[t]
\centering\small
\caption{Directions under test, on five open checkpoints of 7 to 8B parameters. Chat panel: Qwen2.5-7B-Instruct, Qwen3-8B (thinking off), OLMo-2-1124-7B-Instruct. Reasoning panel: DeepSeek-R1-Distill-Qwen-7B, DeepSeek-R1-Distill-Llama-8B, Qwen3-8B (thinking on). D4G denotes LiReFs tested on GSM8K. D10 (test awareness, \citealp{abdelnabi2025testawareness}), preregistered as secondary, was not run.}
\label{tab:directions}
\begin{tabularx}{\linewidth}{@{}l l X l@{}}
\toprule
ID & Direction & Process target $T$ (per item) & Panel \\
\midrule
D1 & Refusal \citep{arditi2024refusal} & the completion refuses (phrase list) & chat \\
D2 & Truth, mass-mean \citep{marks2024geometry} & the model's own true/false verdict & chat \\
D3 & CR axis \citep{ramdan2026crv1} & logit-lens \citep{nostalgebraist2020logitlens} settling depth of the answer & chat \\
D4, D4G & LiReFs \citep{hong2025lirefs} & settling depth on D3's items (D4); chain-of-thought necessity on GSM8K (D4G) & chat \\
D5 & Backtracking \citep{ward2025backtracking} & the model's next sentence backtracks (keyword list) & reasoning \\
D6 & Reasoning strength \citep{sheng2025strength} & the model's own thinking length & reasoning \\
D7 & Known-entity awareness \citep{ferrando2025entities} & the attribute is actually recalled & chat \\
D8 & Sycophancy persona \citep{chen2025persona} & the answer flips to the user's stated wrong option & chat \\
D9 & True thinking \citep{zhao2025truethinking} & the step's True Thinking Score & reasoning \\
\bottomrule
\end{tabularx}
\end{table}

Table~\ref{tab:directions} lists the directions. Each was rebuilt in each panel model with its authors' code, calibration data, layer rule and reading position, adapted where a rule did not port (the adaptations are listed per direction in the Stage~2 documents). Checkpoints that require accepting a licence on a Hugging Face account (Llama, Gemma) were not used, because the runs did not use the author's account (Stage~1, \S0.3); the chat panel replaces those home models for D1, D2, D4 and D7. All models run in bf16 with their chat templates on one RTX~5090.

\paragraph{Preregistration.} Stage~1 (\texttt{PREREGISTRATION.md} in the release: the battery, pass criteria, verdict rules, SESOIs, process targets and hypotheses) was frozen on 14~October~2026 before any code was written. Stage~2 fixed per-direction settings (norm grids, item counts, generation lengths, matched-set constructions, judge choice and Check~0a tolerances) after pilots on Qwen3-1.7B with separate items, and was frozen in parts, D1--D3 on 14~October and D4--D9 on 15~October, each before its main runs (freeze times in Appendix~\ref{app:deviations}). The documents, with the SHA-256 hash and time of each freeze, were kept locally and made public together with the results on 18~October~2026; no third-party registry was used\ifanon.\else{} (\href{https://github.com/Somebodyhere101/direction-validity-battery}{github.com/Somebodyhere101/direction-validity-battery}, file \texttt{PREREG\_SHA256SUMS}).\fi{} Every departure is listed in Appendix~\ref{app:deviations}.

\paragraph{Judges.} Two targets were to be scored by a Qwen3-8B judge: refusal and backtracking. Each judge was validated against blind labels on 50 pilot outputs and was to be used only at Cohen's $\kappa \ge 0.7$. The refusal judge reached 0.52 and the backtracking judge 0.64, so the authors' phrase lists are the primary measures, as Stage~1 requires; judge-based results are reported as secondary.

\section{Results}\label{sec:results}

\subsection{Controls: refusal and truth}

\paragraph{Refusal (D1).} The refusal direction, rebuilt with \citet{arditi2024refusal}'s selection pipeline in each model, predicts which prompts the model actually refuses, including benign prompts that look harmful (AUC 0.92 to 0.99). It beats both label-free directions and moves refusal by 20 to 87 points beyond random pushes of the same output change. In Qwen3-8B the effect of both signs is carried by later layers, and the verdict is \verdict{VALID DRIVER}. In Qwen2.5-7B-Instruct the $+d$ push is carried by later layers, while the $-d$ push acts as an output bias (a push at the last block alone recovers 94\% of it); by Check~D's sign rule the verdict rests on $+d$ and is \verdict{VALID DRIVER}. The positive control passes, so H1 is supported. In OLMo-2 the direction selected at block 23 of 32 acts like an output bias in both signs: re-imposing its clean component at any later block cancels the effect, and a push at the last block alone reproduces 94 to 98\% of it (\verdict{OUTPUT-BIAS}). With the judge's labels instead of phrase matching, every model still tracks (AUC 0.93 to 0.99) and steers in both directions. No vocabulary-matched harmful/harmless set passed the leakage test (XSTest's contrast pairs, \citealp{rottger2024xstest}: bag-of-words AUC 0.75), so the label check is not evaluable for refusal.

\paragraph{Truth (D2).} \citet{marks2024geometry}'s mass-mean direction, built from city statements, tracks the model's own true/false verdict on held-out statements in all three models ($\rho$ 0.63 to 0.80) and beats both label-free directions; neither of these passes Check~A in Qwen2.5-7B-Instruct or Qwen3-8B, while in OLMo-2 $N_\text{pc1}$ does, less strongly ($\rho$ 0.37 [0.23, 0.48]). It separates true from false statements almost perfectly on a matched set (``$X$ is larger than $Y$'' against ``$X$ is smaller than $Y$'' over the same numbers; AUC 0.98 to 1.00; bag-of-words AUC 0.46). Pushing it moves verdicts beyond matched random pushes by 22 and 29 points ($-d$ and $+d$) in OLMo-2 and by 22 points with $-d$ in Qwen2.5-7B-Instruct, where $+d$ moves them by only 7; both are \verdict{VALID DRIVER}, in Qwen2.5-7B-Instruct on the $-d$ sign. In Qwen3-8B a one-gap push moves verdicts by at most 3 points although an oracle at the same layer does steer them: there the direction tracks but does not drive (\verdict{VALID TRACKER}).

\subsection{Known-answer case: the CR axis (D3, H7)}
As in \citet{ramdan2026crv2}, the CR axis does not track settling depth in Qwen2.5-7B-Instruct or Qwen3-8B ($\rho$ $+0.03$ and $+0.07$, both CIs include 0), and pushing it moves settling depth by at most 0.030 of depth (under one layer) beyond matched random pushes. In OLMo-2 it does track ($\rho$ $+0.30$), but $N_\text{len}$ and $N_\text{pc1}$ track better ($+0.43$, $+0.41$). What the push moves is reply format: whether the model shows working, by 21 to 52 points (in Qwen3-8B only with $-d$). The verdict is \verdict{TARGET NOT LINEAR} in all three models, because the oracle's one-gap push does not move settling depth at that layer either. H7 is therefore only partly retrodicted: the tracking and driving failures reproduce \citet{ramdan2026crv2}, but the verdict is not the predicted \verdict{NUISANCE} or \verdict{NOT SHOWN}, and in Qwen2.5-7B-Instruct the push toward ``compute'' raises shown working 14 points more than the length direction does.

\subsection{Process directions}

\paragraph{LiReFs (D4, D4G).} \citet{hong2025lirefs}'s reasoning-versus-memorisation direction was rebuilt with their layer sweep, which scales the component along each layer's candidate direction and keeps the layer that most improves chain-of-thought accuracy on MMLU-Pro reasoning questions \citep{wang2024mmlupro}. The rebuilt direction separates reasoning from recall questions (two-fold AUC 0.97, 0.98 and 0.91 in Qwen2.5-7B-Instruct, Qwen3-8B and OLMo-2), which passes the Check~0a gate of Stage~2b. That gate was redefined after the pilot, which had shown the authors' intervention lowering reasoning accuracy: from the authors' headline effect, Stage~1's definition of Check~0a, to separation alone (a deviation marked $\ast$ in Appendix~\ref{app:deviations}). The intervention at the selected layer does not improve held-out reasoning accuracy in any model (20.7\% to 20.7\%, 42.5\% to 40.2\%, 26.4\% to 19.5\%), so under the Stage~1 definition D4 would be \verdict{NOT REPRODUCED} in all three models. In Qwen3-8B and OLMo-2 the sweep selects block 6, where a weight-randomised copy separates the classes about as well as the trained model (AUC 0.88 against 0.98; 0.92 against 0.91).

Because ``reasoning versus recall'' is the claim of the CR axis too, D4 was tested on D3's items. On the format-matched multiplication items of \citet{ramdan2026crv2}, the direction tracks settling depth weakly in Qwen2.5-7B-Instruct ($\rho$ $+0.19$ [$+0.06$, $+0.32$], short of the SESOI) and less well than $N_\text{len}$ ($+0.31$ [$+0.18$, $+0.43$], close to the $+0.32$ reported there), and not in Qwen3-8B ($-0.08$ [$-0.21$, $+0.04$]) or OLMo-2 ($+0.10$ [$-0.03$, $+0.22$]). A one-gap push moves settling depth by at most 0.015 of depth (0.4 layers) in every model, against the claimed sign in Qwen2.5-7B-Instruct and Qwen3-8B. It moves reply format instead, in opposite directions: the rate of shown working rises by 13 points beyond matched random pushes in Qwen2.5-7B-Instruct (24 beyond label-permuted ones) and falls by 20 points in OLMo-2. A second, independently published reasoning direction thus reproduces the CR axis's pattern: it moves the format of the reply, not the depth of the computation. On GSM8K \citep{cobbe2021gsm8k} (D4G), where chain of thought raises accuracy by 59 to 66 points on average, the projection does not follow how much it helps on each problem in any model ($\rho$ $+0.02$, $-0.04$, $-0.06$), and a push changes direct-answer accuracy by less than the 10-point SESOI, with CIs that exclude it. The oracle tracks settling depth in every model but its one-gap push does not move it; for chain-of-thought necessity the oracle's reading is inconclusive in Qwen2.5-7B-Instruct ($\rho$ $+0.17$ [$+0.03$, $+0.30$]) and fails in Qwen3-8B at block 6. The verdict on both item sets is \verdict{TARGET NOT LINEAR} in all three models.

\paragraph{Backtracking (D5).} \citet{ward2025backtracking}'s backtracking direction, rebuilt from their annotated chains, was tested against whether each sentence of the model's own chain backtracks. With the preregistered measure, the authors' keywords ``Wait'' and ``Hmm'', it does not track (AUC 0.47 to 0.54 over about 11,000 sentences per model), and a one-gap push barely moves the keyword rate. At the authors' own strength, 7.7 gaps, it steers in all three models: $+d$ adds 4.7 to 12.2 backtracking sentences per 1,000 tokens and $-d$ removes 3.3 to 9.7 beyond matched Gaussian pushes (summed KL 33 to 148), against both references except $+d$ in Qwen3-8B. With the judge, which agreed better with blind labels than the keywords did ($\kappa$ 0.64 against 0.54) but missed the preregistered bar, the same vector tracks judged backtracking in all three models (AUC 0.73, 0.69, 0.70; CIs above 0.66) and beats the length direction (0.61 to 0.63). A one-gap push also moves judged backtracking beyond both random references in R1-Distill-Llama-8B and in Qwen3-8B, where $-d$ removes 1.7 judged backtracking sentences per 1,000 tokens [$-2.4$, $-1.1$]; in R1-Distill-Qwen-7B the effect is inconclusive. The two measures agree on 88 to 92\% of sentences, and the disagreements decide the result. By the preregistered rules the verdict is \verdict{TARGET NOT LINEAR} (and \verdict{NOT REPRODUCED} in R1-Distill-Llama-8B, where the rebuilt vector's cosine with the released one was 0.79 against a gate of 0.80); the oracle was pushed only up to two gaps (Section~\ref{sec:strength}). Label-permuted directions are nearly orthogonal to this vector ($|\cos| \le 0.33$); 1 of the 60 tracks the keyword target and none steers it.

\paragraph{Reasoning strength (D6).} \citet{sheng2025strength}'s reasoning-strength direction is a difference of means between MATH \citep{hendrycks2021math} problems of higher and lower difficulty at the last prompt token, which they report the model uses to plan how long to think. It was rebuilt from their calibration file with their layer rule and tested against the model's own thinking length (tokens before \texttt{</think>}, four samples per problem, capped at 4,096). In R1-Distill-Qwen-7B, their home model, their every-layer steering reproduces strongly, and the direction tracks thinking length (partial $\rho$ $+0.45$ [$+0.33$, $+0.55$], controlling for difficulty level, subject and question length). A one-gap $-d$ push shortens thinking by 272 tokens beyond matched random pushes [$-372$, $-178$] against a SESOI of 206, and the effect is handed on to later layers. But both label-free directions track as well ($N_\text{len}$ $+0.47$, $N_\text{pc1}$ $+0.46$), and $N_\text{pc1}$, which is nearly the same vector ($\cos = 0.95$), also steers as much at matched KL ($-286$ tokens [$-383$, $-187$]); $N_\text{len}$ does not steer ($-96$ [$-184$, $-12$]), and $d$ beats it by 176 tokens [98, 255]. The verdict is \verdict{NUISANCE}, the only one in the study. The label survives matching (AUC 0.73 [0.69, 0.77]; bag-of-words AUC 0.63). In R1-Distill-Llama-8B, their steering at the strength they give for that model does not change thinking length by the preregistered 10\% (\verdict{NOT REPRODUCED}), and neither does it in Qwen3-8B with thinking on, so the overall verdict is \verdict{NOT REPRODUCED}. Where the direction does reproduce, it is nearly the first within-class principal component of its own calibration prompts and does no better than it.

\paragraph{Known-entity awareness (D7).} \citet{ferrando2025entities} found sparse-autoencoder latents in Gemma-2 that separate entities the model knows from entities it does not, and showed that steering with the unknown-entity latent makes the model refuse to answer about entities it knows. Gemma is gated, so the direction was rebuilt in the chat panel as the difference-of-means analogue of their latents: per-model known and unknown labels from the model's own accuracy on attribute questions (their labelling rule), averaged within entity type, at the layer with the best minimum-over-types held-out separation. It separates known from unknown entities (minimum over types 0.81 to 0.90), but the refusal effect does not reproduce: pushing the unknown-entity direction on known entities leaves refusal at 8\%, 0\% and 10\% (from 4\%, 4\% and 6\%; random pushes of the same size 14\%, 4\% and 5\%). By the preregistered gate the verdict is \verdict{NOT REPRODUCED} in all three models.

The later checks, reported as specified, are the strongest of any process direction. The projection tracks whether the model actually recalls the asked attribute in every model (AUC 0.68 [0.59, 0.76], 0.66 [0.58, 0.75], 0.77 [0.69, 0.84]) and beats both label-free directions, although 7, 7 and 8 of the 20 label-permuted directions also pass Check~A while lying nearly orthogonal to it ($|\cos| \le 0.36$; Section~\ref{sec:discussion}). The known/unknown label survives matching on pairs matched for token count and entity type (AUC 0.86 to 0.90, bag-of-words 0.42 to 0.53). In Qwen2.5-7B-Instruct a one-gap push raises hallucination on unknown entities by 13 points beyond matched random pushes [6, 21], beats both label-free directions at matched KL, and is handed on to later layers (hand-off at block 16 of 28; a last-block push recovers 34\% of the effect). In Qwen3-8B and OLMo-2 the one-gap push effect is inconclusive (8 and 10 points, CIs from 1 to 20) and the oracle's one-gap push does not pass Check~D; in OLMo-2 a half-gap push of the direction raises hallucination by 13 points [5, 21] against both references (exploratory). Had the gate passed, the analysis code would have given \verdict{VALID DRIVER} in Qwen2.5-7B-Instruct and \verdict{TARGET NOT LINEAR} in the other two, where the oracle's push is inconclusive in one sign and fails in the other; Stage~1's wording, ``fails D'', read literally, would give \verdict{VALID TRACKER} there instead. Our difference-of-means rebuild on these chat models follows the process it is named after and can push the model toward confabulation, but it does not produce the refusal effect that \citet{ferrando2025entities} report for SAE latents in Gemma-2; we did not test their latents or their model.

\paragraph{Sycophancy persona vector (D8).} \citet{chen2025persona}'s sycophancy vector was rebuilt with their extraction recipe (contrastive system prompts, judge-filtered responses, response-averaged difference of means at their relative depth), with a local judge in place of their API judge and ten rather than a hundred samples per question. It reproduces their steering result in every model: adding twice the vector raises judged sycophancy on their evaluation questions from 27 to 91, 57 to 99 and 25 to 88 (coherence 81, 35 and 63 under steering).

The process target asks what sycophancy means for an answer: does the model abandon an answer it gets right without the user's belief for the wrong option the user says they believe? On ARC questions \citep{clark2018arc} with the answer prefilled to a letter, flips are rare in the Qwen models (17 of 1,906 and 42 of 1,931 correctly answered questions) and common in OLMo-2 (508 of 1,540). The projection at the last prompt token tracks flips in Qwen2.5-7B-Instruct (AUC 0.72 [0.62, 0.83] over 14 Check~A flips), and beats both label-free directions there. It does not track them in Qwen3-8B (0.47 [0.35, 0.58]) or in OLMo-2 (0.47 [0.44, 0.50]), where the number of flips leaves little room for doubt. An oracle built from flips tracks them well in every model (AUC 0.98, 0.77, 0.87), so giving in is linearly readable at the vector's layer, just not along the vector. Neither the vector nor the oracle moves the flip rate by more than 2.4 points with a one-gap push in either sign, so the verdict is \verdict{TARGET NOT LINEAR} in all three models. On the matched set (the same belief sentence naming the correct option or a wrong one; bag-of-words AUC 0.37), the vector is higher when the user is right (AUC 0.10, 0.18, 0.36; all \verdict{REVERSED}). It behaves more like agreement with the user than like pressure to give in: the persona vector steers a sycophantic style of reply, which is what its authors' evaluation measures, but it is not shown to carry the decision to give up a correct answer.

\paragraph{True thinking (D9).} \citet{zhao2025truethinking} score each step of a chain of thought by its causal effect on the final answer (the True Thinking Score, TTS) and report a direction that separates steps the model actually uses from decorative ones; pushing it makes perturbed steps decide the answer more or less often. No step labels are released, so each model's direction was rebuilt from its own chains on AMC competition problems (as in the authors' release) with TTS recomputed by their procedure, and tested against the TTS of every numeric step in 240 held-out MATH chains per model. In R1-Distill-Qwen-7B the rebuilt vector's cosine with their released vector is 0.48, just under the preregistered 0.50, although the authors' engagement test passes (+8 points over random vectors), so the verdict there is narrowly \verdict{NOT REPRODUCED}. In R1-Distill-Llama-8B (cosine 0.58) and Qwen3-8B (no released vector) the gate passes. In all three models the direction tracks which steps matter (partial $\rho$ $+0.37$ [$+0.33$, $+0.41$], $+0.29$ [$+0.25$, $+0.34$], $+0.37$ [$+0.33$, $+0.41$]) and beats both label-free directions in two of them; in R1-Distill-Llama-8B the within-class first component tracks as well ($+0.32$). A one-gap push does not change whether the answer follows a perturbed step beyond matched random pushes (R1-Distill-Qwen-7B: $-d$ $-9.5$ points [$-21.0$, $+0.8$], $+d$ $+1.2$ [$-6.2$, $+9.2$]; R1-Distill-Llama-8B: $0.0$ [$-6.7$, $+6.5$] and $-1.7$ [$-8.0$, $+4.1$]), and in Qwen3-8B a $+d$ push makes the answer follow the perturbed step 24 points \emph{less} often, against the claimed sign. The oracle's one-gap push does not pass Check~D in any model, so the verdict is \verdict{TARGET NOT LINEAR} overall. The matched step sets leaked in two models (bag-of-words AUC 0.76 and 0.79); in R1-Distill-Qwen-7B, where the set was accepted, the label survives matching on only 14 pairs (AUC 0.93 [0.81, 1.00]).

\begin{table}[t]
\centering\small
\caption{Verdicts per direction and model, by the preregistered rules. Overall: the verdict in at least two of three models. TNL = target not linear; NR = not reproduced. Home models: R1-Distill-Llama-8B for D5, R1-Distill-Qwen-7B for D6 and D9. Tracks: models in which the direction passes Check~A and beats both label-free directions (Check~C), whatever the verdict; D6's two unreproduced models stopped after calibration.}
\label{tab:verdicts}
\begin{tabular}{@{}l l l l l c@{}}
\toprule
Chat panel & Qwen2.5-7B & Qwen3-8B & OLMo-2-7B & Overall & Tracks \\
\midrule
D1 Refusal & valid driver & valid driver & output bias & valid driver & 3/3 \\
D2 Truth & valid driver & valid tracker & valid driver & valid driver & 3/3 \\
D3 CR axis & TNL & TNL & TNL & TNL & 0/3 \\
D4 LiReFs (D3's items) & TNL$^{a}$ & TNL$^{a}$ & TNL$^{a}$ & TNL & 0/3 \\
D4G LiReFs (GSM8K) & TNL$^{a}$ & TNL$^{a}$ & TNL$^{a}$ & TNL & 0/3 \\
D7 Known-entity awareness & NR$^{b}$ & NR$^{b}$ & NR$^{b}$ & NR & \textbf{3/3} \\
D8 Sycophancy persona & TNL & TNL & TNL & TNL & 1/3 \\
\midrule
Reasoning panel & R1-Qwen-7B & R1-Llama-8B & Qwen3-8B think & Overall & Tracks \\
\midrule
D5 Backtracking & TNL & NR$^{c}$ & TNL & TNL & 0/3 \\
D6 Reasoning strength & nuisance & NR & NR & NR & 0/1 \\
D9 True thinking & NR$^{d}$ & TNL & TNL & TNL & \textbf{2/3} \\
\bottomrule
\end{tabular}

\vspace{0.3em}
\parbox{\linewidth}{\footnotesize $^{a}$Passes a Check~0a gate redefined after the pilot as separation alone; under Stage~1's gate on the authors' headline effect it would be NR. $^{b}$Difference-of-means analogue of SAE latents, on non-home models. $^{c}$Cosine with the released vector 0.79 against a gate of 0.80. $^{d}$Narrowly: cosine with the released vector 0.48 against a gate of 0.50.}
\end{table}

Table~\ref{tab:verdicts} collects the verdicts. Table~\ref{tab:checks} shows what lies behind them for the process directions: the published direction tracks its process in 8 of the 19 process runs that went past calibration (21 in all; D6's two unreproduced models stopped after calibration) and beats both label-free directions in 6, but passes Check~D in only 2, and the oracle, which reads the process in 17 of the 19, passes Check~D in none.

\begin{table}[t]
\centering\small
\caption{Check outcomes behind the verdicts, over the 19 evaluated process-direction runs (of 21: D4 on both item sets and D5--D9, with D6's two unreproduced models stopped after calibration) and the 6 control runs. ``Passes'' as defined in Section~\ref{sec:battery}; Check~D is the one-gap push.}
\label{tab:checks}
\begin{tabular}{@{}l c c@{}}
\toprule
 & Process runs & Control runs \\
\midrule
Direction tracks the process (passes A) & 8 / 19 & 6 / 6 \\
\quad and beats both label-free directions (C) & 6 / 19 & 6 / 6 \\
Direction's push passes D & 2 / 19 & 5 / 6 \\
Oracle reads the process (passes A) & 17 / 19 & 6 / 6 \\
Oracle's push passes D & 0 / 19 & 6 / 6 \\
\bottomrule
\end{tabular}
\end{table}

\subsection{What a null push of the oracle shows}\label{sec:oracle}
The oracle is the difference between the mean activations of the top and bottom terciles of $T$ on 120 oracle items. It is built to read the target, and nothing optimises it to push. A push along it that moves the target shows that the target can be pushed at $\ell^*$; a push that does not move it does not show that no direction can. The data show that this gap is real. In three runs the published direction moves its target, against both references, where the oracle does not. In reasoning strength's home model the direction passes Check~D at one gap, while the oracle's push passes against the Gaussian reference ($-228$ tokens [$-321$, $-133$]) but not against the permuted one, which leaves it inconclusive. In known-entity awareness in Qwen2.5-7B-Instruct the direction passes Check~D, while the oracle's push raises hallucination by only 6 points [1, 12] against the Gaussian reference, short of the SESOI. And in known-entity awareness in OLMo-2, a half-gap push of the direction raises hallucination by 13 points [5, 21] against both references, while the oracle moves it at no strength tested. Against the Gaussian reference alone, as in Section~\ref{sec:strength}, two more runs show the pattern at two gaps: known-entity awareness in Qwen3-8B, and true thinking in R1-Distill-Llama-8B, where the oracle's two-gap cells are not evaluable. The positive controls show only that the same construction can push decisions that are themselves the output. \verdict{TARGET NOT LINEAR} therefore means that the oracle did not read the target or that its one-gap push did not pass Check~D. It does not mean that the target cannot be pushed linearly at the published layer, and a published direction's failure under that verdict is uninformative only in the weak sense that the battery cannot tell it apart from a failure of the oracle's push. Figure~\ref{fig:oracle} shows the oracle's reading and its one-gap push for every run.

\begin{figure}[t]
\centering
\begin{tikzpicture}
\begin{axis}[
  width=0.82\linewidth, height=6.4cm,
  xmin=-0.3, xmax=5.3, ymin=-1.6, ymax=3.25,
  xlabel={Oracle's reading: Check~A statistic in SESOI units},
  ylabel={Oracle's one-gap push (SESOI units)},
  axis lines=left, tick align=outside, tick style={black!50},
  axis line style={black!60},
  label style={font=\small}, tick label style={font=\footnotesize},
  legend style={font=\footnotesize, draw=none, fill=white, fill opacity=0.85, text opacity=1, at={(0.99,0.02)}, anchor=south east, row sep=-1pt},
  legend cell align=left,
  clip=false,
]
\addplot[black!35, dashed, forget plot] coordinates {(-0.3,1) (5.3,1)};
\addplot[black!35, dashed, forget plot] coordinates {(1,-1.6) (1,3.25)};
\addplot[black!20, forget plot] coordinates {(-0.3,0) (5.3,0)};
\addplot[only marks, mark=*, mark size=3.2pt, color=kdec, draw=white, line width=0.5pt] coordinates {(4.66,3.00) (4.59,3.00) (4.14,3.00) (4.15,3.00) (3.99,3.00) (4.65,3.00)};
\addlegendentry{decision (D1, D2)}
\addplot[only marks, mark=square*, mark size=2.8pt, color=kans, draw=white, line width=0.5pt, error bars/.cd, y dir=both, y explicit, error bar style={line width=0.6pt, color=kans}, error mark=none] coordinates {(2.98,0.80) += (0,0.98) -= (0,0.96) (2.15,0.63) += (0,0.59) -= (0,0.50) (2.52,0.36) += (0,0.74) -= (0,0.68) (3.67,0.22) += (0,0.50) -= (0,0.58) (4.79,0.24) += (0,0.25) -= (0,0.18) (2.72,0.08) += (0,0.19) -= (0,0.15) (2.05,-0.68) += (0,0.70) -= (0,0.76) (1.97,0.14) += (0,0.79) -= (0,0.73) (2.17,0.67) += (0,0.76) -= (0,0.68)};
\addlegendentry{answer choice (D7, D8, D9)}
\addplot[only marks, mark=triangle*, mark size=3.6pt, color=kcomp, draw=white, line width=0.5pt, error bars/.cd, y dir=both, y explicit, error bar style={line width=0.6pt, color=kcomp}, error mark=none] coordinates {(2.69,-0.05) += (0,0.31) -= (0,0.29) (2.47,0.04) += (0,0.04) -= (0,0.04) (1.10,0.59) += (0,0.13) -= (0,0.13) (2.74,-0.07) += (0,0.10) -= (0,0.11) (2.21,0.40) += (0,0.07) -= (0,0.07) (1.26,-0.01) += (0,0.02) -= (0,0.02) (1.36,0.10) += (0,0.31) -= (0,0.29) (0.84,0.29) += (0,0.55) -= (0,0.53) (0.14,-0.01) += (0,0.09) -= (0,0.12) (1.84,-0.04) += (0,0.38) -= (0,0.40) (2.18,0.28) += (0,0.43) -= (0,0.41) (1.90,0.32) += (0,0.31) -= (0,0.32) (2.36,1.11) += (0,0.45) -= (0,0.46)};
\addlegendentry{course of computation (D3--D6)}
\node[font=\footnotesize, text=black!70, anchor=east] at (axis cs:3.9,3.0) {6 control runs, 4.2 to 8.7 (off scale) $\uparrow$};
\node[font=\footnotesize, text=black!70, anchor=west] at (axis cs:2.45,1.18) {D6};
\node[font=\footnotesize, text=black!55, anchor=south east] at (axis cs:5.3,1.0) {SESOI};
\end{axis}
\end{tikzpicture}
\caption{Every run's oracle, read and pushed. Horizontal: the oracle's Check~A statistic in SESOI units ($\rho/0.20$, or $(\text{AUC}-0.5)/0.10$ for binary targets; 1 is the pass threshold). Vertical: its one-gap push against the matched-KL Gaussian reference, in SESOI units, for the sign with the larger effect in the claimed direction, with 95\% CI. The six control runs lie off scale and are drawn at the top edge without intervals. Among process runs only reasoning strength's home model (D6) reaches the SESOI, and against the permuted reference it is inconclusive.}
\label{fig:oracle}
\end{figure}

\subsection{Exploratory: push strength}\label{sec:strength}
The preregistered pushes are one calibration gap. Two analyses ask what other strengths do; neither enters any verdict. First, every main run also pushed $d$, the label-free directions and the oracle at 0.5 and 2 gaps, strengths that Stage~1 lists as exploratory for Check~D; we computed their matched-KL effects with the preregistered analysis code (D6's oracle was pushed at one gap only). Second, a sweep that is post hoc and not in Stage~1: it was added on 18~October~2026, after every verdict was final, and pushed $d$ and the oracle at 4 and 8 gaps for D3, D4, D4G and D8, on the first 80 Check~D items, with the run's own oracle and Gaussian directions and the random grid extended to 24 gaps. The sweep script covers only these four directions; D5, D6, D7 and D9 were not swept, and no other reason was recorded. A target counts as moved at a strength when, in a cell with at least half of its items evaluable, the effect against matched Gaussian pushes reaches the SESOI with its CI excluding zero in the push's sign, and the other sign is not a significant reversal (Check~D's sign rule; the permuted reference was not rerun). An item is not evaluable when its pushed output cannot be scored or the push's KL falls outside the random grid; Appendix~\ref{app:evaluable} gives the evaluable count of every cell that lost items. For backtracking we also report the published direction at its authors' strength (7.7 gaps).

\begin{table}[t]
\centering\footnotesize
\setlength{\tabcolsep}{4.2pt}
\caption{Exploratory: the smallest push, in calibration gaps, at which the oracle ($o$) or the published direction ($d$) moves its target beyond matched-KL Gaussian pushes. ``--\,($k$)'': no movement at any evaluable strength up to $k$ gaps. Strengths 0.5 to 2 come from the main runs; 4 and 8 from the post hoc sweep (D3, D4, D4G, D8 only). n.r.: not run (stopped after calibration).}
\label{tab:strength}
\begin{tabular}{@{}l l cc cc cc@{}}
\toprule
 & & \multicolumn{2}{c}{Qwen2.5-7B} & \multicolumn{2}{c}{Qwen3-8B} & \multicolumn{2}{c}{OLMo-2-7B} \\
\cmidrule(lr){3-4}\cmidrule(lr){5-6}\cmidrule(l){7-8}
Target (direction) & Kind & $o$ & $d$ & $o$ & $d$ & $o$ & $d$ \\
\midrule
Refusal (D1) & decision & 0.5 & 0.5 & 0.5 & 0.5 & 0.5 & 0.5 \\
True/false verdict (D2) & decision & 0.5 & 1 & 0.5 & 2 & 0.5 & 0.5 \\
Hallucination, unknown entities (D7) & answer choice & 2 & 1 & --\,(2) & 2 & --\,(2) & 0.5 \\
Flip to the user's option (D8) & answer choice & 2 & --\,(8) & 4 & --\,(8) & 8 & --\,(8) \\
Settling depth (D3) & computation & --\,(8) & --\,(8) & 4 & 4 & --\,(2)$^{b}$ & --\,(4)$^{c}$ \\
Settling depth (D4) & computation & --\,(4)$^{d}$ & --\,(8) & --\,(8) & --\,(8) & --\,(8) & --\,(8) \\
Chain-of-thought necessity (D4G) & computation & --\,(8) & --\,(8) & --\,(8) & --\,(8) & --\,(8) & --\,(8) \\
\midrule
 & & \multicolumn{2}{c}{R1-Qwen-7B} & \multicolumn{2}{c}{R1-Llama-8B} & \multicolumn{2}{c}{Qwen3-8B think} \\
\cmidrule(lr){3-4}\cmidrule(lr){5-6}\cmidrule(l){7-8}
Answer follows a perturbed step (D9) & answer choice & 2 & 2 & --\,(1)$^{e}$ & 2 & --\,(2) & --\,(2) \\
Backtracking, keywords (D5) & computation & --\,(2) & 7.7$^{a}$ & --\,(2) & 7.7$^{a}$ & --\,(2) & 2 \\
Thinking length (D6) & computation & 1$^{f}$ & 1 & n.r. & n.r. & n.r. & n.r. \\
\bottomrule
\end{tabular}

\vspace{0.3em}
\parbox{\linewidth}{\footnotesize $^{a}$At the authors' strength; no movement at 2 gaps or less. $^{b}$Replies broke down above 2 gaps; the sweep is not evaluable. $^{c}$At 2 gaps $-d$ lowers settling depth beyond the SESOI, but $+d$ lowers it too (a significant reversal), so by the sign rule it does not count; above 2 gaps only $+d$ at 4 gaps is evaluable. $^{d}$Not evaluable at 8 gaps (1 of 80 items in one sign, none in the other); at 4 gaps evaluable in one sign only. $^{e}$Both signs not evaluable at 2 gaps (27 of 61 and 15 of 59 items); the $+$ cell passes on its 15 items. $^{f}$Pushed at one gap only; passes against the Gaussian reference but not the permuted one.}
\end{table}

Table~\ref{tab:strength} separates the targets only in part. Decisions that are themselves the output (refuse or not; true or false) move along the oracle at half a gap in all six runs. Answer choices move along the oracle in 5 of 9 runs, at 2 to 8 gaps: hallucination in Qwen2.5-7B-Instruct, giving in in all three chat models (for example by 59 points of flips at four gaps in Qwen2.5-7B-Instruct [47, 70]), and following a perturbed step in R1-Distill-Qwen-7B. The four other runs were evaluable only up to two gaps, or one gap for true thinking in R1-Distill-Llama-8B. Settling depth and chain-of-thought necessity moved along the oracle in 1 of the 8 runs evaluable at four or eight gaps: Qwen3-8B's settling depth at four gaps, by $-0.055$ of depth (about 2 of 36 layers). D3 in OLMo-2, whose replies broke down above two gaps, is not counted. Strength in gaps is not strength in output change. The strongest evaluable oracle pushes on these two targets changed the output by a median summed KL of 13 (D4 in OLMo-2) to 198 (D4G in Qwen2.5-7B-Instruct) without moving them, and in six of the seven runs by more than the largest KL, 27, at which an answer choice moved (3 to 27 over the five runs that moved). Backtracking did not move along the oracle, which was pushed only to two gaps (median summed KL at most 6), and thinking length moved at one gap in its one run.

The published directions moved their own targets in 9 of the 19 process runs at some strength: known-entity awareness in all three models (0.5 to 2 gaps), true thinking in two (2 gaps), backtracking in all three at its authors' strength (in Qwen3-8B already at 2 gaps, by $+1.9$ keyword sentences per 1,000 tokens [$+1.1$, $+2.8$] against a SESOI of 1.4), and reasoning strength in its home model (1 gap). The CR axis moved settling depth once (Qwen3-8B, $-0.033$ of depth, about 1.2 layers, at four gaps). LiReFs never moved settling depth or chain-of-thought necessity, and the sycophancy vector never moved giving in, including at the strengths at which the oracle does. The gradient that survives is narrow: along the oracle, answer choices move at strengths and output changes at which settling depth and chain-of-thought necessity do not.

\subsection{Hypotheses}
\begin{table}[t]
\centering\small
\caption{Preregistered hypotheses and outcomes.}
\label{tab:hyp}
\begin{tabularx}{\linewidth}{@{}l X X@{}}
\toprule
 & Prediction & Outcome \\
\midrule
H1 & Refusal is \verdict{VALID DRIVER} in at least two of three models. & Supported (Qwen2.5-7B-Instruct, on $+d$; Qwen3-8B; output bias in OLMo-2). \\
H2 & Label-permuted directions fail A and D in every model. & Not supported. Where the labelled contrast dominates the calibration variance they keep much of the real direction ($|\cos|$ up to 0.95) and many pass A; only D5's are clean nulls. No verdict changes if Check~D uses the Gaussian reference alone (Section~\ref{sec:discussion}). \\
H3 & For at least half the directions, a weight-randomised model separates held-out calibration prompts at AUC $\ge 0.80$. & Not supported: 4 of 9 (refusal, CR axis, LiReFs, sycophancy). \\
H4 & Norm-matched effects exceed KL-matched effects. & Supported: $+0.25$ [$+0.07$, $+0.49$] SESOI over 23 reproduced pairs; median $+0.01$, 16 of 23 positive, and two pairs give $+2.0$ and $+1.6$. \\
H5 & Fewer than half the process directions are valid overall. & Supported: 0 of 6. \\
H6 & At least two process directions receive \verdict{NUISANCE} in some model. & Not supported: 1 (D6 in R1-Distill-Qwen-7B). \\
H7 & The battery retrodicts \citet{ramdan2026crv2} for the CR axis. & Partly: tracking and driving fail as there, but the verdict is \verdict{TARGET NOT LINEAR}, not the predicted \verdict{NUISANCE} or \verdict{NOT SHOWN}, and in Qwen2.5-7B-Instruct the axis beats $N_\text{len}$ on shown working. \\
\bottomrule
\end{tabularx}
\end{table}

Table~\ref{tab:hyp} gives the outcomes. H6's prediction fails partly because of the verdict order: the length or first-component direction matches or beats the published direction at tracking in four further runs (the CR axis in two models, LiReFs in Qwen2.5-7B-Instruct, true thinking in R1-Distill-Llama-8B), but each of those verdicts is pre-empted by \verdict{TARGET NOT LINEAR}.

\section{Discussion}\label{sec:discussion}

\paragraph{Readable at the published layer, pushable only in part.} The most common verdict is \verdict{TARGET NOT LINEAR}: at the preregistered strength the oracle reads 17 of 19 process targets at the published layer and its one-gap push passes Check~D for none of them, while for refusal and truth it passes in all six runs. The dissociation itself is not new: \citet{xun2026lagged} reports it for probe directions across the Pythia suite and its training, and other studies for hallucination, hazard, truth and reasoning-correctness probes (Section~\ref{sec:related}). Our oracle results add per-item targets measured from what the model does, for processes that published directions claim to track, and they come with a caveat. The same construction pushes refusal and truth at half a gap, but those are decisions that are themselves the output, so this shows only that a tercile-mean oracle can push such decisions; it does not show that the oracle's null on a process target is a property of the target rather than of the oracle, and in at least three runs the published direction pushes where the oracle does not (Section~\ref{sec:oracle}). \citet{xun2026lagged}'s pushes add the unit-norm probe direction on a fixed strength grid, while ours are scaled by the calibration gap, so strengths are not directly comparable across the two studies. The exploratory strengths (Table~\ref{tab:strength}) leave a narrower gradient than a wall. Along the oracle, answer choices (whether to hallucinate, to give in, to follow a perturbed step) moved in 5 of 9 runs at summed KLs of 3 to 27, while settling depth and chain-of-thought necessity moved in 1 of 8 runs, although in six of the other seven the oracle's strongest push changed the output more than that. Backtracking and thinking length do not fit a simple reading in which the course of a computation stays put: the published backtracking vector moves the keyword rate at its authors' strength in all three models (the oracle was never pushed that far), and thinking length, which is also a surface property of the output, moves at one gap along both directions. For settling depth and chain-of-thought necessity, a plausible reading is that one additive push at one layer can bias what the model says without changing how deep the answer forms or whether written reasoning helps; those properties may be carried by non-linear or distributed features that a single push does not reach. This fits recent accounts of where steering works: where a concept is aligned with the unembedding \citep{billa2026lap}, at depths that differ from where probes read best \citep{wang2026readbest}, and along directions organised by the model's distribution over answers \citep{yu2026answerbasin}. A decision that is itself the output, or a choice among answers, has such an alignment; settling depth need not, and reasoning-correctness directions in a small trained model likewise fail to improve reasoning under intervention \citep{nwadike2026correctness}. Either way, a demonstration that a direction ``steers reasoning'' at one strength and layer is not evidence that the process it reads is linearly controllable there, and a null push along a reading direction is not evidence that it is not.

\paragraph{Tracking and reproduction come apart.} Known-entity awareness, which fails its reproduction gate in every model, and true thinking, which fails it only in its home model and narrowly (cosine with the released vector 0.48 against 0.50), are the best trackers in the study. Known-entity awareness follows recall on the model's own items and beats both label-free directions in all three models; true thinking follows its steps' True Thinking Score in all three and beats both in two; and their labels survive matching wherever a matched set was accepted. For known-entity awareness the headline that fails is a steering effect (refusal on known entities), tested on a difference-of-means rebuild in non-home models; this matches \citet{galeone2026detection}, who find the direction that detects fake entities about $83^\circ$ from the one that produces refusal in Gemma~2-2B-it. For true thinking the gate fails only on the cosine with the released vector, while the authors' engagement test passes. Both directions do push their own targets in five of six runs at 0.5 to 2 gaps (Table~\ref{tab:strength}), effects that the one-gap Check~D found in only one. The converse also holds: refusal and truth, whose headline is steering, steer. Reproduction of a steering demonstration and validity as a tracker are separate questions, and the battery shows they come apart in both directions.

\paragraph{What the push moves instead.} Where a reasoning direction moves anything, it often moves the form of the reply. LiReFs, like the CR axis before it, changes whether the model shows working in Qwen2.5-7B-Instruct and OLMo-2, in opposite signs, while leaving settling depth unchanged. The sycophancy persona vector reliably makes replies sycophantic in style, which is what its authors measure, but it is higher when the user is right than when the user is wrong and does not change whether the model gives up a correct answer. Reasoning strength, where it reproduces, is nearly the within-class first principal component of its calibration prompts. In each case a style, format or length variable that co-varies with the calibration contrast is a better description than the process in the direction's name.

\paragraph{Label-permuted directions are not a clean null.} Stage~1 called H2 an implementation check whose failure ``means a bug, not a finding''. We do not read the failure as a bug, for two reasons, although we have not otherwise audited the permutation code. First, where the labelled contrast dominates the calibration activations' variance, as for refusal, truth, the CR axis and the sycophancy vector, a random relabelling keeps much of the real direction ($|\cos|$ up to 0.95), and many permuted directions track the outcome; this is SteerCheck's alignment leakage \citep{luo2026steercheck} in a second setting. Second, where the contrast is small relative to other variance, the same code gives clean nulls: for backtracking ($|\cos| \le 0.33$) 1 of 60 permuted directions passes Check~A and none steers. One case fits neither: for known-entity awareness the permuted directions are nearly orthogonal to $d$ ($|\cos| \le 0.36$), yet 7 to 8 of 20 pass Check~A in each model. They are pooled relabellings, while $d$ is averaged within entity type, and the difference of means of a random split lies mostly along the activations' largest-variance directions, which here may separate entity types that differ in how often the attribute is recalled; we have not tested this. Requiring Check~D to beat the permuted reference changed no verdict: with the Gaussian reference alone every verdict is the same. KL-matched Gaussian directions are the more reliable reference. Matching them on KL rather than norm matters because candidate feature directions move the output more than isotropic random ones of the same size \citep{ferreiradasilva2026error}. Here norm-matched references overstated effects by a quarter of a SESOI on average (H4), but the median excess was $+0.01$ SESOI and 16 of 23 pairs were positive; two pairs, refusal in Qwen2.5-7B-Instruct ($+2.0$) and true thinking in Qwen3-8B ($+1.6$), carry most of the mean.

\paragraph{The measurement of the process decides the verdict.} For backtracking, two measures that agree on about nine sentences in ten give opposite tracking verdicts: the authors' keywords count ``Hmm, let me break this down'' and miss re-checks without those words. For sycophancy, flips are so rare in the Qwen models that tracking rests on 14 and 36 positive items, and the model with many flips gives the clearest answer. A process target is a measurement, and its validity, base rate and power have to be established like the direction's.

\paragraph{Limitations.} The oracle is a difference of tercile means on 120 items, a reading direction that nothing optimises for pushing; it is not an upper bound on pushability, and in at least three runs the published direction pushes where it does not, so \verdict{TARGET NOT LINEAR} means only that the oracle's push was null at one gap. The exploratory strength analyses cover 0.5 to 2 gaps for every run but 4 and 8 gaps only for four directions, were added after the verdicts (the 4- and 8-gap sweep is post hoc), use the Gaussian reference only, lose items in some cells (Appendix~\ref{app:evaluable}), and cover no other layers. Some verdicts rest on one push sign, as the sign rule allows: refusal in Qwen2.5-7B-Instruct on $+d$, truth in Qwen2.5-7B-Instruct and reasoning strength on $-d$, and known-entity awareness in Qwen2.5-7B-Instruct on $+d$. Both preregistered judges missed their validation bar, so refusal and backtracking rest on phrase lists that agreed only moderately with blind labels in the pilot. The panel is five checkpoints of 7 to 8B parameters. Home models that need a licence agreement were not used: D1, D2, D4 and D7 were tested on a chat panel rather than the models they were published on, and D7 only in a difference-of-means form, which may explain part of its reproduction failure. LiReFs' reproduction gate was redefined after its pilot. Several layer rules were adapted (D2, D7, D9), D6 and D9 cap generation length, and D8 uses a local judge and fewer extraction samples. One GPU set the budget: D6 uses the preregistered reduced design in two models, and D10 was not run. The verdicts are per model without correction across models or directions. The breadth is across published directions and models: other studies of the read/push gap go further across training and scale \citep{xun2026lagged} and across layers \citep{billa2026lap,wang2026readbest}.

\section{Conclusion}
A preregistered battery that passes the refusal and truth controls finds no valid process direction among six published ones. The commonest verdict is that the oracle reads the target at the published layer but its one-gap push does not move it, which is a statement about that push and not proof that the target cannot be pushed. At larger exploratory strengths, answer choices move along the oracle in more runs than settling depth or chain-of-thought necessity do, and the published directions move their own targets in 9 of 19 process runs, backtracking among them at its authors' strength. Behind the verdicts, known-entity awareness and true thinking track their processes well while failing all or part of their reproduction gates, and several directions that do steer move style, format or length variables that ride along with the calibration contrast. Calibration separation, which weight-randomised models reproduce for four of nine directions, and fixed-norm steering demonstrations, which norm-matched baselines flatter, cannot tell these cases apart. Measuring the process per item, comparing with label-free directions from the same prompts, matching pushes on KL, and building an oracle from the target can, within the limits set out above. The code, preregistrations and per-run results are released so that the battery can be applied to other directions.

\paragraph{Code and data.} \ifanon Links to the code, preregistrations and results are withheld for anonymous review.\else Code, preregistrations with their SHA-256 freeze records, the deviation log and per-run results: \url{https://github.com/Somebodyhere101/direction-validity-battery}; archived with the per-item results at \href{https://doi.org/10.5281/zenodo.23496934}{doi:10.5281/zenodo.23496934}.\fi{} No direction vectors are released.

\paragraph{Use of AI assistance.} The idea of packaging the checks from the author's earlier studies as a battery and applying it to published directions is the author's. An AI system (Claude, Anthropic), working under the author's direction, reviewed the literature, wrote the preregistrations and code, ran the experiments and analyses, and drafted this paper. The preregistrations record that the author had not reviewed their details line by line when they were frozen. \textcolor{red}{\textbf{[AUTHOR-ROLE-TODO: the author's own sentence on what the author personally did and checked goes here. It must be written by the author.]}} The author directed the work and is responsible for its content.

\bibliographystyle{plainnat}
\bibliography{references}

\appendix
\section{Deviations from the preregistration}\label{app:deviations}
Every departure from Stage~1 is listed here with its reason; the full text, with the numbering used in the release, is in the Stage~2 documents and \texttt{DEVIATIONS\_LOG.md} in the release. Departures that change what is measured are marked $\ast$; the rest change cost, engineering or batching only. Stage~2 freeze times (AST): D1--D3 on 14~October; on 15~October, D5 at 01:03, D4 at 12:12, D8 at 16:19, D7 at 17:53, D6 at 19:15 and D9 at 19:35.

\paragraph{All directions.}
\begin{itemize}
\item $\ast$ Checkpoints that need a licence agreement (Llama, Gemma) were not used; the chat panel replaces them (Stage~1, \S0.3, fixed before any code).
\item $\ast$ Both preregistered judges failed validation (refusal $\kappa$ 0.52, backtracking 0.64), so the authors' phrase lists are primary, as Stage~1 specifies; judge results are secondary.
\item D10 (test awareness), preregistered as secondary, was not run; by Stage~1 it then enters neither H5 nor H6.
\item A performance-only patch was applied to the frozen Stage~2a code midway through D1, after OLMo-2's attention cache spilled into system memory. It changes floating-point reduction order only; OLMo-2's D1 push was restarted so that every item used one version.
\item Pushes that ran out of GPU memory on long prompts (D4G in all three chat models, D8 in Qwen3-8B) were restarted from the first item in smaller row chunks (96, or 48 for OLMo-2), so that batch size stays fixed within each model's comparison. Partial files are kept and not analysed.
\item Two output fields were added to the analysis summaries after the Stage~2b-D5 freeze; nothing computed changed.
\item The 4- and 8-gap strength sweep (Section~\ref{sec:strength}) was added on 18~October, after every verdict was final. It is post hoc: Stage~1 lists 0.5 and 2 gaps as exploratory strengths, but not this sweep. The deviation log's entry for it cites Stage~1 \S9, which does not list it.
\item Stage~1's format clamp, and with it \verdict{FORMAT-MEDIATED}, was never evaluated, because no target was measured both free and under a prefill.
\end{itemize}

\paragraph{D1--D3 (Stage 2a).} The backtracking pilot ran after, not before, the D1--D3 main runs, and on the RTX~5090 rather than the laptop. $\ast$ D2's layer rule is a substitute (best held-out mass-mean probe accuracy), because Marks \& Tegmark chose layers by patching. Arditi et al.'s selection pipeline was reimplemented rather than run from their repository. No matched set exists for D3 (Check~B not run); D1's was rejected for leakage.

\paragraph{D4.} One fixed 200-question sample is used for every layer of the sweep (theirs draws anew per layer); the five-shot exemplars were refetched from MMLU-Pro's validation split. $\ast$ The GSM8K drive target is direct-answer accuracy, because pushing 384-token sampled chains on every reference row was not affordable. $\ast$ The Check~0a gate is on separation only, decided after the pilot showed their intervention lowering reasoning accuracy; the intervention's effect is reported. Under Stage~1's gate on the authors' headline effect, D4 would be \verdict{NOT REPRODUCED} in all three models.

\paragraph{D5.} $\ast$ Ward et al.'s recipe and annotations are used in every panel model, where Stage~1 named Venhoff et al.'s recipe (which needs GPT-4o annotations) for the non-home models. The sentence splitter and the sentence-start token mapping were fixed in the pilot. Check~B is not constructed. The secondary judge pass over pushed outputs was first limited to the home model to free the GPU and then completed.

\paragraph{D6.} $\ast$ Thinking is capped at 4,096 tokens (theirs 16,384); capped chains count at the cap. $\ast$ Random and permuted reference rows use one sample each rather than the four-sample mean, for cost. The preregistered long-generation reduction ($R = 8$, 5 norms, 80 items) is applied to R1-Distill-Llama-8B and Qwen3-8B. The layer is their code's single-layer index, not the layer named in their text. Sampling uses top-$p$ 0.95 as in their generation scripts. Models that fail the gate skip the later stages. D6 uses its own sampled-generation engine.

\paragraph{D7.} $\ast$ The direction is the difference-of-means analogue of their SAE latents, averaged within entity type, with a substitute layer rule (best minimum-over-types held-out separation). Calibration uses a random 600 entities per type, labelled per model. $\ast$ The oracle is built within entity type, because the pooled oracle mostly encoded entity type. Check~D items come from a held-out pool selected per model by label, so that each model has known and unknown items; the drive target combines refusal on known and hallucination on unknown entities into one per-sign target. The label-permuted directions are pooled relabellings, not averaged within entity type.

\paragraph{D8.} $\ast$ Ten samples per extraction question (theirs 100), 400 new tokens (theirs 1,000), and a local Qwen3-8B judge with their prompts (theirs gpt-4.1-mini); the judge is used only for the extraction filter and the Check~0a gate. $\ast$ The answer is prefilled with ``The answer is ('' after two free-answer pilots left up to a quarter of answers without a letter. $\ast$ The matched set names the correct or a wrong option in the same belief sentence, because ``with and without the belief'' cannot pass the leakage rule. Check~A was extended to every unused ARC item before the main runs.

\paragraph{D9.} $\ast$ Every model's direction is rebuilt from its own chains with True Thinking Scores recomputed by the authors' procedure (no step labels are released); two bugs in their scoring code (a dropped delimiter, a double BOS token) are fixed, the decorative cut-off is $\mathrm{TTS} \le 0.005$ as in their paper, and confidence is teacher-forced. Chains are capped (8,192 tokens for calibration, 4,096 for test). Test problems are MATH Levels 1--3 with integer answers; Check~A uses numeric steps only. The layer is their released example's (17 of 28, same relative depth elsewhere). The engagement and disengagement tests used in the gate were written afresh, since the release does not include them.

\section{Evaluable items in the exploratory push cells}\label{app:evaluable}
Table~\ref{tab:evaluable} gives, for every run in Table~\ref{tab:strength}, the items pushed per sign and every cell in which fewer were evaluable. An item is not evaluable when its pushed output cannot be scored or the push's summed KL falls outside the random grid. A cell is a kind of push ($o$ or $d$), a strength in gaps and a sign; cells marked $\ast$ have fewer than half their items evaluable and are treated as not evaluable. All other cells, including D5 at its authors' strength, are complete. The sweep's denominator is the first 80 Check~D items (for D8, those of them with a scorable unpushed outcome).

\begin{table}[H]
\centering\scriptsize
\caption{Items per sign ($-$/$+$ where they differ) and cells with fewer evaluable items.}
\label{tab:evaluable}
\begin{tabularx}{\linewidth}{@{}l l c c X@{}}
\toprule
Dir. & Model & Main & Sweep & Cells with fewer evaluable items (kind, gaps, sign: evaluable/items) \\
\midrule
D1 & Qwen2.5-7B & 47/73 & -- & all complete \\
D1 & Qwen3-8B & 47/73 & -- & all complete \\
D1 & OLMo-2-7B & 47/73 & -- & all complete \\
D2 & Qwen2.5-7B & 120 & -- & $o$\,1$+$ 98/120; $o$\,2$+$ 92/120; $d$\,2$+$ 91/120 \\
D2 & Qwen3-8B & 120 & -- & $o$\,2$+$ 113/120 \\
D2 & OLMo-2-7B & 120 & -- & all complete \\
D3 & Qwen2.5-7B & 120 & 80 & $o$\,2$+$ 116/120; $d$\,2$+$ 108/120; $o$\,8$-$ 49/80; $o$\,8$+$ 79/80; $d$\,8$-$ 0/80$^\ast$ \\
D3 & Qwen3-8B & 120 & 80 & $o$\,2$+$ 118/120; $o$\,8$-$ 26/80$^\ast$; $o$\,8$+$ 12/80$^\ast$; $d$\,8$+$ 0/80$^\ast$ \\
D3 & OLMo-2-7B & 120 & 80 & $o$\,4$-$ 0/80$^\ast$; $o$\,4$+$ 0/80$^\ast$; $d$\,4$-$ 0/80$^\ast$; $o$\,8$-$ 0/80$^\ast$; $o$\,8$+$ 0/80$^\ast$; $d$\,8$-$ 0/80$^\ast$; $d$\,8$+$ 0/80$^\ast$ \\
D4 & Qwen2.5-7B & 120 & 80 & $o$\,2$-$ 114/120; $o$\,2$+$ 113/120; $d$\,2$-$ 100/120; $o$\,4$-$ 37/80$^\ast$; $o$\,4$+$ 63/80; $d$\,4$-$ 2/80$^\ast$; $o$\,8$-$ 1/80$^\ast$; $o$\,8$+$ 0/80$^\ast$; $d$\,8$-$ 0/80$^\ast$ \\
D4 & Qwen3-8B & 120 & 80 & $o$\,2$-$ 118/120; $o$\,8$-$ 43/80; $o$\,8$+$ 72/80; $d$\,8$+$ 71/80 \\
D4 & OLMo-2-7B & 120 & 80 & $o$\,0.5$+$ 119/120; $d$\,2$+$ 99/120; $d$\,4$-$ 3/80$^\ast$; $d$\,8$-$ 0/80$^\ast$ \\
D4G & Qwen2.5-7B & 120 & 80 & $o$\,2$-$ 78/120; $o$\,2$+$ 119/120; $d$\,2$-$ 102/120 \\
D4G & Qwen3-8B & 120 & 80 & all complete \\
D4G & OLMo-2-7B & 120 & 80 & $d$\,2$-$ 96/120; $d$\,8$+$ 78/80 \\
D7 & Qwen2.5-7B & 60 & -- & $d$\,2$+$ 59/60 \\
D7 & Qwen3-8B & 60 & -- & $d$\,2$+$ 54/60 \\
D7 & OLMo-2-7B & 60 & -- & $o$\,2$+$ 59/60; $d$\,2$+$ 52/60 \\
D8 & Qwen2.5-7B & 183 & 73 & $d$\,0.5$-$ 180/183; $o$\,2$+$ 182/183; $d$\,2$+$ 178/183; $d$\,8$+$ 16/73$^\ast$ \\
D8 & Qwen3-8B & 190 & 74 & $o$\,0.5$+$ 188/190; $d$\,0.5$-$ 188/190; $d$\,0.5$+$ 187/190; $d$\,1$+$ 189/190; $o$\,2$-$ 151/190; $o$\,2$+$ 186/190; $d$\,8$+$ 38/74 \\
D8 & OLMo-2-7B & 165 & 69 & $o$\,1$-$ 159/165; $o$\,2$-$ 21/165$^\ast$; $o$\,2$+$ 96/165; $d$\,2$-$ 160/165 \\
D5 & Qwen3-8B & 120 & -- & all complete \\
D5 & R1-Qwen-7B & 120 & -- & all complete \\
D5 & R1-Llama-8B & 120 & -- & all complete \\
D6 & R1-Qwen-7B & 120 & -- & all complete \\
D9 & Qwen3-8B & 50/70 & -- & $o$\,2$-$ 49/50; $o$\,2$+$ 38/70; $d$\,2$+$ 67/70 \\
D9 & R1-Qwen-7B & 58/62 & -- & $o$\,2$+$ 59/62 \\
D9 & R1-Llama-8B & 61/59 & -- & $o$\,2$-$ 27/61$^\ast$; $o$\,2$+$ 15/59$^\ast$; $d$\,2$-$ 60/61; $d$\,2$+$ 49/59 \\
\bottomrule
\end{tabularx}
\end{table}

\end{document}
